% GitHub Repo and Documentation: https://github.com/celiobjunior/resume-template
% Copyright © 2025 Celio B Junior. All rights reserved.
% 
% Licensed under the Apache License, Version 2.0 (the "License");
% you may not use this file except in compliance with the License.
% You may obtain a copy of the License at
%
%     http://www.apache.org/licenses/LICENSE-2.0
%
% This template follows best practices from README.md

% Início do documento LaTeX: define tipo e formato do documento
% a4paper = tamanho A4, 10pt = tamanho base da fonte
\documentclass[a4paper,10pt]{article}

% --- PACOTES: BASE / IDIOMA ---
\usepackage[utf8]{inputenc}
\usepackage[T1]{fontenc}
\usepackage[portuguese]{babel}

% --- PACOTES: LAYOUT / TIPOGRAFIA ---
\usepackage{geometry}
\usepackage{parskip}
\usepackage{microtype}
\usepackage{enumitem}
\usepackage{titlesec}
\usepackage{array}
\usepackage{tabularx}

% --- PACOTES: LINKS / BOOKMARKS ---
\usepackage{hyperref}
\usepackage{bookmark}
\usepackage{xurl}

% --- CONFIGURAÇÃO: MARGENS / ESPAÇAMENTO ---
\geometry{top=1.5cm, bottom=1.5cm, left=1.5cm, right=1.5cm} % Margens da página
\setcounter{secnumdepth}{0}                                 % Remove numeração de seção
\setlist[itemize]{
    leftmargin=0.75em, % Recuo da lista (mais perto da margem da página)
    itemsep=0.2em,     % Espaço entre itens
    topsep=0.25em,     % Espaço antes/depois da lista
    parsep=0em,        % Espaço entre parágrafos no mesmo item
    partopsep=0em      % Espaço extra ao iniciar a lista
}

% Remove números de página e cabeçalhos para visual limpo do currículo
\pagestyle{empty}

% --- CONFIGURAÇÃO: METADADOS / LINKS ---
\hypersetup{
    pdftitle={CV Lucas Ribeiro de Lima},   % Título do PDF
    pdfauthor={Lucas Ribeiro de Lima},     % Autor do PDF
    colorlinks=true,          % Usa cor nos links (sem caixa)
    linkcolor=black,          % Cor de links internos
    urlcolor=black,           % Cor de URLs
    citecolor=black,          % Cor de citações
    bookmarksdepth=2          % Barra lateral: seção e subseção
}

% --- CONFIGURAÇÃO: TÍTULOS ---
\titleformat{\section}
{\Large\bfseries}
{}
{0em}
{}
[\titlerule\vspace{0.5ex}]
\titleformat{\subsection}
{\normalsize\bfseries}
{}
{0em}
{}
\titlespacing*{\subsection}{0pt}{0.35em}{0.15em}

% Macro para entradas do currículo com bookmark no PDF
\newcounter{cventry}
\newcommand{\cventry}[4]{%
    \refstepcounter{cventry}%
    \phantomsection%
    \pdfbookmark[2]{#1}{cventry-\thecventry}%
    \noindent\begin{tabularx}{\textwidth}{@{}>{\raggedright\arraybackslash}X >{\raggedleft\arraybackslash}X@{}}
    \textbf{#1} & #2 \\
    \textit{#3} & \textit{#4} \\
    \end{tabularx}
}

% --- INÍCIO DO DOCUMENTO ---
\begin{document}

% --- CABEÇALHO ---
\begin{center}
    {\LARGE \textbf{Lucas Ribeiro de Lima}} 
    \\ [0.1cm]
    Guarulhos, São Paulo, Brasil
    {\textbullet}
    \href{mailto:lucasribeiro292004@hotmail.com}{lucasribeiro292004@hotmail.com} 
    {\textbullet}
    \href{https://www.linkedin.com/in/lucas-ribeiro-de-lima-2b83b3231/}{linkedin.com/in/lucas-ribeiro-de-lima-2b83b3231} 
    {\textbullet}
    \href{https://github.com/kasslima}{github.com/kasslima}
\end{center}

% --- SEÇÕES ---

\section{Resumo}
    Engenheiro de Software com foco em desenvolvimento Full Stack e experiência na construção de aplicações web, APIs, automações de processos, soluções de dados e integrações entre sistemas. Atuo com Node.js, React, C\#/.NET, PostgreSQL, Python, Azure Functions, AWS e pipelines CI/CD, participando da implementação, deploy, sustentação em produção e melhoria contínua. Atualmente aprofundo conhecimentos em Go, arquitetura cloud, sistemas distribuídos e desenho de soluções escaláveis.

\section{Experiência}
    \cventry{Recurso Digital}{São Paulo, SP}{Desenvolvedor Full Stack Pleno}{Abr 2026 - Atual}
        \begin{itemize}
            \item Desenvolvo e mantenho aplicações web em todo o ciclo de desenvolvimento, da definição técnica à sustentação em produção.
            \item Defino soluções técnicas para novas funcionalidades, evolução de sistemas, ganhos de escalabilidade e modernização de legados.
            \item Construo APIs, integrações e serviços backend com Node.js, C\#/.NET e Azure Functions.
            \item Automatizei processos de negócio com Python e n8n, reduzindo fluxos manuais de aproximadamente 8 horas para poucos minutos.
            \item Estruturo e mantenho pipelines CI/CD para automatizar builds, testes, deploys e atualizações de ambiente.
            \item Modelo e otimizo consultas em PostgreSQL para aplicações, análises, dashboards e relatórios operacionais.
        \end{itemize}

    \cventry{Recurso Digital}{São Paulo, SP}{Desenvolvedor Full Stack Junior}{Out 2025 - Abr 2026}
        \begin{itemize}
            \item Desenvolvi e mantive aplicações web full stack, contribuindo com serviços backend e funcionalidades frontend.
            \item Implementei APIs, serviços e integrações entre sistemas usando Node.js e C\#/.NET.
            \item Construí interfaces em React e mantive Azure Functions em ambientes Microsoft Azure.
            \item Criei automações em Python para reduzir tarefas repetitivas e melhorar a eficiência de processos internos.
            \item Atuei na sustentação de sistemas em produção, incluindo deploys, monitoramento, análise de incidentes e manutenção de pipelines CI/CD.
        \end{itemize}

    \cventry{Recurso Digital}{São Paulo, SP}{Estagiário em Análise de Dados}{Jul 2025 - Out 2025}
        \begin{itemize}
            \item Desenvolvi soluções de dados, automações de processos e funcionalidades web para otimizar fluxos operacionais.
            \item Construí rotinas de ETL para coleta, transformação, validação e carregamento de dados.
            \item Criei dashboards e análises em Power BI para acompanhar indicadores e apoiar a tomada de decisão.
            \item Contribuí com funcionalidades frontend e backend usando React, Node.js, C\#/.NET, PostgreSQL, Azure e Azure Functions.
            \item Usei Git para apoiar versionamento, colaboração e entrega ao longo do ciclo de desenvolvimento.
        \end{itemize}

    \cventry{Foundever}{São Paulo, Brasil}{Analista de Suporte Jr}{Jun 2024 - Jul 2025}
        \begin{itemize}
            \item Prestei suporte técnico para problemas de hardware e software em ambientes de clientes.
            \item Gerenciei chamados por canais digitais, garantindo comunicação clara com clientes e acompanhamento de incidentes.
            \item Usei Salesforce e Sprinklr para registrar, acompanhar, priorizar e resolver casos de suporte.
            \item Monitorei indicadores de atendimento e desempenho operacional usando Microsoft Excel.
            \item Colaborei com equipes técnicas para escalonar, investigar e resolver incidentes complexos.
        \end{itemize}

% ------

\section{Habilidades}
    \begin{itemize}
        \item \textbf{Linguagens e Frameworks:} Node.js, React, C\#/.NET, Python, Go, JavaScript, REST APIs
        \item \textbf{Dados e Automação:} PostgreSQL, Power BI, ETL, n8n, dashboards, relatórios, automação de processos
        \item \textbf{Cloud e DevOps:} AWS, Amazon EC2, Azure Functions, CI/CD, Git, balanceamento de carga
        \item \textbf{Idiomas:} Português (Nativo), Inglês (Profissional Completo)
    \end{itemize}

% ------

\section{Certificações}
    \begin{itemize}
        \item \textbf{AWS Technical Essentials} -- Amazon Web Services (AWS), Jun 2026
        \item \textbf{AWS Cloud Practitioner - Training Badge} -- Amazon Web Services (AWS), Jun 2026
        \item \textbf{AWS Cloud Practitioner Essentials} -- Amazon Web Services (AWS), Jun 2026
        \item \textbf{Derive Insights from BigQuery Data Skill Badge} -- Google, Jun 2025
        \item \textbf{Microsoft Power BI Para Business Intelligence e Data Science} -- Data Science Academy, Jan 2025
        \item \textbf{Fundamentos de Linguagem Python Para Análise de Dados e Data Science} -- Data Science Academy, Out 2023
        \item \textbf{Curso de Excel do Básico ao Avançado} -- Udemy, Set 2023
        \item \textbf{Curso de Python} -- Udemy, Abr 2023
        \item \textbf{Algoritmos e Lógica de Programação} -- Udemy, Abr 2021
        \item \textbf{Inglês} -- Kumon Brasil, Jan 2018
        \item \textbf{Pacote Office} -- Microlins, Fev 2018
    \end{itemize}

% ------

% !!!!!!!!!!!!!!!!!!!!!!
% Mude esta seção para o começo, depois do cabeçalho e antes
% das suas experiências se você está procurando sua primeira
% vaga ou algum estágio. Do contrário, deixe como está.
% !!!!!!!!!!!!!!!!!!!!!!

\section{Educação}
    \cventry{Centro Universitário Eniac}{Guarulhos, SP}{Ensino Superior, Tecnologia da Informação}{Jan 2024 - Jun 2026}

\end{document}
